\section{Problemática}
\label{sec:problematica}

En abril de 2018 el TSB migró la plataforma y cerca de 1,9 millones de clientes se quedaron sin
servicio. El informe de la Cámara de los Comunes cifra el episodio en cientos de millones de libras
y en semanas de operación restringida (House of Commons Treasury Committee, 2018). No habla de
OpenAPI ni de BIAN. Documenta otra cosa: el costo de cambiar el sistema cuando negocio y tecnología
no tienen una correspondencia que se pueda revisar. Retrabajo, adaptadores armados a prueba y error,
daño de reputación.

El \textit{open banking} británico es un caso distinto. La orden de la CMA (2016) obligó, a partir
de 2018, a nueve bancos de \textit{retail} a exponer APIs. Publicar el Swagger no dice cómo se
compara esa interfaz con un modelo de negocio de referencia. Esa comparación sigue en la gobernanza
del banco, no en el archivo.

En lo que leímos, el hueco se repite con otros números. Casas et al.\ (2021) mapean 47 estudios
primarios y solo el 40\,\% ancla un dominio; el resto trabaja APIs genéricas. Golmohammadi et al.\
(2024) revisan 92 trabajos de prueba REST: validar un servicio ya desplegado sale caro, y los casos
industriales que se pueden repetir son pocos. Farzi (2021) acerca BIAN a la agilidad del banco y
deja abierta la distancia entre el \textit{to-be} y lo que corre.

En un contrato el desajuste no es un ensayo. Arquitectura guarda el \ServiceDomain{}; desarrollo
versiona el YAML del \textit{sprint}. Para la misma orden aparecen \texttt{PaymentOrder} y
\texttt{paymentRequest}. Un \Behavior{} declarado con PUT no tiene el GET en el contrato, o al
revés, y la versión nueva del archivo no dispara una revisión del modelo. Cuando hay revisión, es
lectura humana o un \textit{schema matcher} que no conoce el SD. Shabbir et al.\ (2025) atan ese
desfase a la consistencia y a la mantenibilidad de la documentación. El
§\ref{sec:planteamiento} lo formula como problema tecnológico; el Anexo~C, como fila de la matriz.
